% !PRINTED 556
\weberpage{556}

Verfährt man so mit allen Primzahlen und setzt
\[
S'_{10} + S'_{01} - 2\,S'_{11} = T',
\]
so folgt
\[
T' = T\,\Pi\left(1 - \frac{1}{p^2}\right),
\]
worin sich das Produkt $\Pi$ auf \gesperrt{alle ungeraden Primzahlen} $p$ erstreckt. Nun erhält man
durch Entwickelung nach steigenden Potenzen von $p^{-2}$:
\[
\begin{aligned}
\frac{1}{\Pi(1-p^{-2})} &= 1 + \frac{1}{3^2} + \frac{1}{5^2} + \frac{1}{7^2} + \frac{1}{9^2} + \cdots,\\
&= \frac{\pi^2}{8},
\end{aligned}
\]
und folglich
\[
T = \frac{\pi^2}{8}\,T',
\]
und indem man ebenso mit den beiden anderen Summen verfährt, folgt aus (9):
\weq{10}{\begin{aligned}
\log f(\omega_1)\,f(\omega_2) &= \tfrac{1}{3}\log 2 + \frac{\beta\pi}{24}\left(S'_{10} + S'_{01} - 2\,S'_{11}\right), \\
\log f_1(\omega_1)\,f_1(\omega_2) &= \tfrac{1}{3}\log 2 + \frac{\beta\pi}{24}\left(S'_{10} - 2\,S'_{01} + S'_{11}\right), \\
\log f_2(\omega_1)\,f_2(\omega_2) &= \tfrac{1}{3}\log 2 + \frac{\beta\pi}{24}\left(-2\,S'_{10} + S'_{01} + S'_{11}\right).
\end{aligned}}

Diese Formeln vereinfachen sich wesentlich, wenn man $\alpha = 0$, also
\[
\omega_1 = \omega_2 = \omega = i\,\beta
\]
setzt, also wenn man ein \gesperrt{rein imaginäres} $\omega$ annimmt. Dann wird
\[
S'_{10} = \sum \frac{1}{x^2 - \omega^2 y^2} \qquad \begin{array}{l}x \text{ ungerade}, \ y \text{ gerade},\\xy \text{ relativ prim}.\end{array}
\]

Sondert man das Glied $x = \pm 1$, $y = 0$ ab und nimmt von den übrigen je vier zusammen,
so ergibt sich
\[
S'_{10} = 2 + 4\sum \frac{1}{x^2 - \omega^2 y^2},
\]
worin $x$, $y$ nur \gesperrt{positive Werte} durchlaufen.

Ebenso verfährt man mit den anderen Summen. Setzt man dann
%CONT>>
